\section{Joint forest selection in the abstract occurrence model}
\label{sec:joint}

The top-handle theorem allows many possible deletion trees. The rows left
behind can differ, so it is natural to ask whether a tree can be chosen to
make the subsequent one/two-handle cancellation easy. This section solves
an exact abstract version of that problem. The geometric constraint in the
next section is indispensable when interpreting the answer.

\subsection{Raw occurrences, saturation, and peeling}

Fix the primal forest used to cancel the zero/one-handle forest. The
remaining one-handles have labels $C$, and the original two-handles have
labels $E$. Let
\[
 a_{ec}\in\mathbb Z_{\geq0}
\]
be the actual unsigned number of passages of the attaching curve of $e$
through the cocore of $c$. These are geometric occurrences, not freely
reduced words or signed exponent sums. For the native core they are obtained
by counting the dimension-one flags of the source attaching cycles after
omitting the primal forest columns.

\begin{definition}[Unit-column peeling]
For a live row set $R\subseteq E$, a pivot $(e,c)$ is a unit-column pivot
if $e\in R$ and
\[
 \sum_{f\in R} a_{fc}=a_{ec}=1.
\]
Peeling this pivot deletes row $e$. The other rows remain unchanged, and
column $c$ becomes zero. A row set is peelable if all its rows can be
deleted this way.
\end{definition}

Put $b_{ec}=\min(2,a_{ec})$. Every unit-column predicate on every row
subset is unchanged by this entrywise saturation: a nonnegative sum is one
exactly when it consists of one unit contribution and otherwise zeros. We
therefore use $b$ for peeling, retaining $a$ when integral or mod-two chain
information is needed.

A nonempty row set $Z$ is a \emph{stopping set} if each of its column sums
is zero or at least two. Its first row can never be peeled, even with other
rows present. Conversely, a stalled peeling process leaves a stopping set.
Unions of stopping sets are stopping sets, so their largest union is the
canonical residual. These statements and the equivalence with complete
positive raw cancellation are prior results in the incoming
\emph{Stopping Cores} report~\cite{incoming-stopping}. The present paper
uses them within that exact positive model. Allowing geometric slides with
intersection cancellation or changing the cocore system is a different
problem.

No arbitrary columns may be discarded. If private bundle columns are
omitted, an independent geometric projection theorem must already justify
the omission, as in the saturated-bundle work. The combinatorial theorems
below can be used on any designated matrix, but a manifold interpretation
requires an authenticated producer for precisely that matrix.

\subsection{Short rows form a graph}

Call $e$ short if $\sum_c b_{ec}\leq2$. Introduce an auxiliary ground
vertex $\omega$, distinct from the cap-graph root $\rho$. The occurrence
multigraph $H$ has vertices $C\cup\{\omega\}$ and one labelled edge for
every short row:
\begin{center}
\begin{tabular}{@{}ll@{}}
\toprule
Nonzero entries of the short row & Its edge in $H$\\
\midrule
One unit at $c$ & $c\omega$\\
Two units at distinct $c,d$ & $cd$\\
One entry at least two at $c$ & A loop at $c$\\
No entries & A loop at $\omega$\\
\bottomrule
\end{tabular}
\end{center}
The last case is necessary: an empty two-handle row cannot be deleted by a
one/two-handle pivot. Parallel edges and loops are genuine cycles.

\begin{lemma}[Graphic peeling criterion]
\label{lem:graphic}
A set of short rows $R$ is completely peelable if and only if $H[R]$ is a
forest.
\end{lemma}

\begin{proof}
At a nonground vertex, the graph degree equals the saturated column sum,
with a loop counted twice. A nonempty forest component has at least two
leaves. At most one leaf is $\omega$, so there is a nonground degree-one
vertex and therefore a legal unit-column pivot. Removing its edge preserves
the forest, and induction peels every row.

Conversely, the edges of any cycle form a stopping set. At every nonground
vertex their degree is zero or at least two. This includes an interior
loop, a loop at $\omega$, and a two-edge cycle of parallel edges. None of
these edge sets can have a first peeled row.
\end{proof}

\subsection{The two-forest criterion}

First suppose every row is short, and let $G$ be the connected cap graph on
the same edge label set $E$. By \cref{lem:rooted}, a complete top schedule
deletes a spanning tree $F$ of $G$. Its complement must be an $H$-forest if
the residual rows are to peel.

\begin{theorem}[Abstract joint forest criterion]
\label{thm:joint}
For connected $G$ and short-row graph $H$, the following are equivalent:
\begin{enumerate}
\item There is a spanning tree $F$ of $G$ such that $H[E\setminus F]$ is
a forest.
\item $E$ is a disjoint union of a $G$-forest and an $H$-forest.
\item For every $Y\subseteq E$,
\begin{equation}
 |Y|\leq r_G(Y)+r_H(Y).
 \label{eq:rank-test}
\end{equation}
\end{enumerate}
A deterministic polynomial-time procedure finds the two schedules or a
subset violating \eqref{eq:rank-test}.
\end{theorem}

\begin{proof}
The first condition immediately gives the second. Given a partition into
forests $I_G,I_H$, extend $I_G$ to a spanning tree by moving edges from
$I_H$. If $I_G$ does not yet span, connectedness of $G$ supplies an edge of
$I_H$ between different $I_G$ components. Adding it preserves the
$G$-forest; removing it preserves the $H$-forest. Repeat. This proves the
first condition and gives the two schedules by \cref{lem:rooted,lem:graphic}.

Intersecting a forest partition with $Y$ proves necessity of
\eqref{eq:rank-test}. Sufficiency is the two-matroid partition theorem of
Edmonds--Fulkerson~\cite[Theorem 1c]{edmondsfulkerson}. For completeness,
the explicit augmenting proof and its witness are given below.
\end{proof}

Now let
\[
 D=\{e:\textstyle\sum_c b_{ec}>2\},\qquad E_0=E\setminus D.
\]
Consider the restricted strategy that requires all rows of $D$ to be
deleted by the top-handle forest. This is a restriction of the strategy,
not a necessary condition for raw cancellation. The single row $(1,1,1)$
is already peelable.

\begin{corollary}[Forced long rows]
\label{cor:forced}
The restricted strategy is possible exactly when $G$ is connected,
$G[D]$ is a forest, and for all $Y\subseteq E_0$,
\begin{equation}
 |Y|\leq r_G(Y\cup D)-|D|+r_H(Y).
 \label{eq:forced-rank}
\end{equation}
Here $H$ is the short-row graph on $E_0$. When $D$ is independent, the first
two terms on the right are $r_{G/D}(Y)$.
\end{corollary}

\begin{proof}
A spanning tree can contain $D$ only when $D$ is a forest. Contracting this
forest identifies spanning trees containing $D$ with spanning trees of
$G/D$. Graphic matroid contraction has rank
$r_{G/D}(Y)=r_G(Y\cup D)-r_G(D)$, and $r_G(D)=|D|$.
Apply \cref{thm:joint}. Contraction can create loops and parallel edges;
they remain in the instance.
\end{proof}

\subsection{A constructive proof with a failure certificate}
\label{sec:augment}

Write $G'=G/D$, with $D=\varnothing$ in the unforced case. Maintain
disjoint sets $I_0,I_1\subseteq E_0$, independent in $G'$ and $H$,
respectively. All other edges are unassigned. An edge $e\notin I_i$ is an
$i$-sink if adding it to $I_i$ preserves a forest. Otherwise its endpoints
are joined by the unique path in the forest $I_i$. Draw an exchange arc
$e\to f$ for each path edge $f$. For a loop this path has no edges and
there is no exchange arc.

Run breadth-first search in this exchange digraph, starting at every
unassigned edge. If a shortest source-to-sink path
$e_0,e_1,\ldots,e_\ell$ is found, put its terminal edge into its sink
forest. For every $j<\ell$, assign $e_j$ to the old forest of
$e_{j+1}$. The number of assigned edges increases by one.

Here is an explicit independence argument. Represent graph edges by their
endpoint incidence vectors over $\F$. Forest edge vectors are independent.
For an incoming edge $e_j$, its expansion in its target forest has
coefficient one on $e_{j+1}$. It has coefficient zero on every later
removed edge $e_{k+1}$ of that same forest, $k>j$: otherwise there would be
an exchange arc skipping part of the shortest path. The matrix of these
coefficients on the removed edges is triangular with diagonal one.
Replacing the removed forest edges by the corresponding incoming edges
therefore preserves the old span and independence. No earlier path vertex
was a sink, since that would give a shorter path. The terminal edge lies
outside the old span of its sink forest and increases its rank by one.
This proves the augmentation rule.

If no sink is reachable, let $Y$ be the reached edge set and $U$ the
unassigned set. Then $U\subseteq Y$. Every fundamental path reached in
either forest lies entirely in $Y$. Consequently
\[
 r_i(Y)=|I_i\cap Y|\quad(i=0,1),\qquad
 |Y|-r_0(Y)-r_1(Y)=|U|.
\]
When $U\ne\varnothing$ this is a violated rank inequality. It also proves
optimality of the current partial partition: any two-forest union has size
at most
\[
 |E_0\setminus Y|+r_0(Y)+r_1(Y)=|E_0|-|U|,
\]
and the current partition attains this value. When $U=\varnothing$, all
edges are assigned and the desired full partition has been found.

For $m_0=|E_0|$ and $v=|V(G')|+|V(H)|$, the simple implementation has at
most $m_0$ augmentations. Each search visits at most $m_0$ edge labels and
uses ordinary forest-path traversals. Its conservative structural bound is
\[
 O\bigl(m_0^2(v+m_0)+m_0^3\log(m_0+2)\bigr).
\]
Integer parsing, label comparisons, and source authentication are additional
polynomial bit costs. Large occurrence values are compared with two, never
expanded into repeated crossings by this planner.

\subsection{An exact residual-deficiency theorem}

When the restricted strategy cannot finish, its failure need not be reported
only as a boolean. Assume $G$ is connected and $D$ is a forest. Define
\begin{equation}
 d=\max_{Y\subseteq E_0}
 \bigl(|Y|-r_{G/D}(Y)-r_H(Y)\bigr).
 \label{eq:deficiency}
\end{equation}
The empty set makes $d\geq0$.

\begin{theorem}[Optimal residual nullity]
\label{thm:deficiency}
The deficiency satisfies
\begin{equation}
 d=\min_{\substack{F\text{ spanning tree of }G\\D\subseteq F}}
 \bigl(|E\setminus F|-r_H(E\setminus F)\bigr).
 \label{eq:minmax}
\end{equation}
The exchange procedure produces an attaining tree, an occurrence forest in
its residue, and a rank witness certifying the common optimum.
\end{theorem}

\begin{proof}
The augmentation argument proves that the maximum two-forest union in
$E_0$ has size $m_0-d$. For any tree $F\supseteq D$, let $J$ be a maximum
$H$-forest in its residue $R=E\setminus F$. Then $F\setminus D$ and $J$
are a two-forest union, so
\[
 |F\setminus D|+r_H(R)\leq m_0-d.
\]
Since $|F\setminus D|+|R|=m_0$, the residual nullity is at least $d$.

For the reverse inequality, take a maximum partial partition $I_G,I_H$.
Every unassigned edge has endpoints in one $I_G$ component; otherwise it
could be assigned to that forest. Hence $I_G\cup I_H$ spans the connected
graph $G/D$. Extend $I_G$ to a spanning tree $T$ using edges of $I_H$ only.
The unused part $J\subseteq I_H$ remains an $H$-forest. Its complement in
$E_0\setminus T$ consists of exactly the $d$ unassigned edges. Therefore
\[
 |E_0\setminus T|-r_H(E_0\setminus T)
 \leq |E_0\setminus T|-|J|=d.
\]
Lift $T$ through the contracted forest $D$. This attains the lower bound.
\end{proof}

The number $d$ is the minimum number of \emph{additional arbitrary row
deletions} that would make an optimized residue a forest. Such deletions are
not licensed geometric cancellations and are never performed by the
prototype. The certificate measures the exact abstract obstacle; it is not
a proof that the manifold is not a handlebody.

\subsection{Two instructive limitations}

\begin{example}[An abstract exchange gain with no geometric realization]
\label{ex:abstract}
Take two interior cap vertices $u,v$ and three rows:
\[
\begin{array}{c|c|c}
e&\text{cap edge in }G&\text{single-column occurrence row}\\\hline
0&uv&(1)\\
1&v\rho&(0)\\
2&u\rho&(2).
\end{array}
\]
Deleting edges $0,1$ leaves the unpeelable row $(2)$; deleting $1,2$
leaves $(1)$ and succeeds. This is a genuine test of the abstract exchange
algorithm. It is not an example of a three-manifold improvement. Over
$\F$, its occurrence boundary and cap incidence have nonzero product:
each interior vertex column sees the unit contribution from row $0$.
It violates $\partial_2\partial_3=0$. \Cref{sec:chain} explains the
consequence of imposing that identity.
\end{example}

\begin{proposition}[General peelability is not a matroid]
\label{prop:not-matroid}
The peelable row subsets of an arbitrary nonnegative occurrence matrix
need not satisfy the matroid augmentation axiom.
\end{proposition}

\begin{proof}
Use the four binary rows
\[
 a=(0,0,1),\quad b=(0,1,1),\quad
 c=(1,1,0),\quad f=(1,1,1).
\]
Both $I=\{c,f\}$ and $J=\{a,b,c\}$ peel: in $I$ first use the third
column of $f$; in $J$ first use the first column of $c$, then the second
of $b$, then the third of $a$. Yet $|J|>|I|$ and neither $I\cup\{a\}$
nor $I\cup\{b\}$ peels. Every nonzero column sum in either enlarged set
is at least two. Thus augmentation fails.
\end{proof}

This counterexample does not prove hardness of the general geometric
problem. It identifies the precise reason that the short-row matroid
argument cannot simply be extended to arbitrary long rows.
